\documentclass[12pt]{article}
\usepackage[spanish]{babel}
\usepackage[utf8]{inputenc}
\usepackage[T1]{fontenc}
\usepackage{amsmath, amssymb}
\usepackage{geometry}
\usepackage{enumitem}
\usepackage{needspace}
\usepackage{tikz}
\usetikzlibrary{babel}

\geometry{letterpaper, margin=2cm}
\setlength{\parindent}{0pt}
\setlength{\parskip}{6pt}

\newcommand{\puntografico}{0.09}

\begin{document}

\begin{center}
  {\Large \textbf{Borrador de items}}\\
  {\large Bloque 1. Geometria - Afirmacion 7}\\[4pt]
  Prueba Nacional Estandarizada de Matematica, Secundaria 2026
\end{center}

\section*{Afirmacion}
Resuelve problemas relacionados con transformaciones en el plano.

\section*{Items propuestos}

\begin{enumerate}[label=\textbf{\arabic*.}, itemsep=1.05cm]

\Needspace{0.95\textheight}
\item
Considere la siguiente representacion grafica para responder los items 1 y 2.

En el plano de diseno de una institucion educativa, el poligono $ABCD$ representa la forma de una zona verde que sera reproducida en diferentes posiciones del patio.

\begin{center}
\begin{tikzpicture}[scale=0.48]
  \def\puntotransformacion{0.21}
  \path[use as bounding box] (-7.9,-1.2) rectangle (1.8,7.8);
  \draw[step=1, gray!45, very thin] (-7,0) grid (0,7);
  \draw[<->, thick] (-7.6,0) -- (1.3,0);
  \draw[<->, thick] (0,-0.7) -- (0,7.4);
  \node[anchor=west] at (1.45,-0.15) {$x$};
  \node[anchor=south east] at (0.65,7.55) {$y$};
  \fill[gray!15] (-2,6) -- (-4,1) -- (-6,4) -- (-5,5) -- cycle;
  \draw[densely dashed, gray!70, thin] (-2,0) -- (-2,6) -- (0,6);
  \draw[densely dashed, gray!70, thin] (-4,0) -- (-4,1) -- (0,1);
  \draw[densely dashed, gray!70, thin] (-6,0) -- (-6,4) -- (0,4);
  \draw[densely dashed, gray!70, thin] (-5,0) -- (-5,5) -- (0,5);
  \draw[very thick] (-2,6) -- (-4,1) -- (-6,4) -- (-5,5) -- cycle;
  \foreach \x/\lab in {-7/{\llap{-}7},-6/{\llap{-}6},-5/{\llap{-}5},-4/{\llap{-}4},-3/{\llap{-}3},-2/{\llap{-}2},-1/{\llap{-}1}}
    \node[anchor=north, inner sep=1pt] at (\x,0) {$\lab$};
  \foreach \y in {1,2,3,4,5,6,7}
    \node[anchor=east, inner sep=1pt] at (0,\y) {$\y$};
  \fill (-2,6) circle (\puntotransformacion) node[anchor=south west, inner sep=1pt, xshift=1pt, yshift=1pt] {\small A};
  \fill (-4,1) circle (\puntotransformacion) node[anchor=north east, inner sep=1pt, xshift=-1pt, yshift=-1pt] {\small B};
  \fill (-6,4) circle (\puntotransformacion) node[anchor=east, inner sep=1pt, xshift=-1pt] {\small C};
  \fill (-5,5) circle (\puntotransformacion) node[anchor=south east, inner sep=1pt, xshift=-1pt, yshift=1pt] {\small D};
\end{tikzpicture}
\end{center}

Si al poligono $ABCD$ se le aplica una reflexion respecto a la recta dada por $y=-x$, entonces, \textquestiondown cual es la imagen del punto $A$?

\begin{enumerate}[label=\Alph*)]
  \item $(-2,-6)$
  \item $(6,-2)$
  \item $(-6,2)$
  \item $(2,6)$
\end{enumerate}

\Needspace{0.52\textheight}
\item
En el mismo plano de diseno, al poligono $ABCD$ se le aplica una homotecia de razon $k=-2$ con centro en el origen y, luego, una traslacion de $3$ unidades hacia la derecha.

\textquestiondown Cual es la imagen del punto $C$ luego de aplicar ambas transformaciones?

\begin{enumerate}[label=\Alph*)]
  \item $(-15,8)$
  \item $(15,-8)$
  \item $(9,-8)$
  \item $(12,-5)$
\end{enumerate}

\Needspace{0.48\textheight}
\item
En una aplicacion de mapas, la ubicacion de una camara de seguridad se representa por el punto $P(3,-2)$ en un plano cartesiano. Para ajustar la orientacion del mapa en la pantalla, el sistema aplica una rotacion de $90^\circ$ en sentido antihorario alrededor del origen.

\textquestiondown Cuales son las coordenadas del punto que representa la nueva ubicacion de la camara en la pantalla?

\begin{enumerate}[label=\Alph*)]
  \item $(-2,-3)$
  \item $(2,3)$
  \item $(-3,2)$
  \item $(3,2)$
\end{enumerate}

\end{enumerate}

\section*{Clave y justificacion breve}

\begin{enumerate}[label=\textbf{\arabic*.}]
  \item \textbf{C.} La reflexion respecto a la recta $y=-x$ transforma un punto $(x,y)$ en $(-y,-x)$. Por tanto, $A(-2,6)$ se transforma en $(-6,2)$.
  \item \textbf{B.} El punto $C(-6,4)$, bajo la homotecia de razon $-2$, se transforma en $(12,-8)$. Al trasladarlo $3$ unidades hacia la derecha se obtiene $(15,-8)$.
  \item \textbf{B.} Una rotacion de $90^\circ$ en sentido antihorario alrededor del origen transforma $(x,y)$ en $(-y,x)$. Por tanto, $P(3,-2)$ se transforma en $(2,3)$.
\end{enumerate}

\end{document}
